\subsection{SPACES OF MAPPINGS INTO A METRIC SPACE}

Let X be a set, Y a uniform space, \((f_{t})_{t\in I}\) a family of pseudometrics defining the uniform structure of Y (Chapter IX, § 1, no. 4), and let S be a set of subsets of X. For each \(t\in I\) , each set \(A\in S\) , and each pair \((u,v)\) of mappings of X into Y, write

\[
g _ {t, \mathrm{A}} (u, v) = \sup _ {x \in \mathrm{A}} f _ {t} (u (x), v (x));
\]

it follows immediately that \(g_{t,\mathbf{A}}\) is a pseudometric on \(\mathcal{F}(\mathbf{X};\mathbf{Y})\) and that the family of pseudometrics \((g_{t,\mathbf{A}})_{t\in \mathbf{I},\mathbf{A}\in \mathfrak{S}}\) defines the uniformity of \(\mathfrak{S}\) -convergence on \(\mathcal{F}(\mathbf{X};\mathbf{Y})\) . In particular:

PROPOSITION I. If Y is a metrizable uniform space, the uniformity of uniform convergence on \(\mathcal{F}(X;Y)\) is metrizable.

For if d is a metric on Y compatible with its uniform structure, the structure of uniform convergence on \(\mathcal{F}(X;Y)\) is defined by the single pseudometric

\[
\delta (u, v) = \sup _ {x \in \mathbf {X}} d (u (x), v (x));
\]

in general this pseudometric is not finite, but it is equivalent to a finite one (Chapter IX, § 1, no. 2), and since the uniformity of uniform convergence is Hausdorff (§ 1, no. 2, Proposition 1), it is metrizable.

COROLLARY. Let X be a topological space and let Y be a metrizable uniform space. Suppose that there is a sequence \((\mathbf{K}_{n})\) of compact subsets of X such that every compact subset of X is contained in some \(K_{n}\) . Then the uniformity of compact convergence on \(\mathcal{F}(\mathbf{X};\mathbf{Y})\) is metrizable.

Since the \(\mathbf{K}_n\) cover \(\mathbf{X}\) , \(\mathcal{F}_c(\mathbf{X};\mathbf{Y})\) is isomorphic to a uniform subspace of \(\prod_{n}\mathcal{F}_{u}(\mathbf{K}_{n};\mathbf{Y})\) ( \(\S 1\) , no. 2, Remark 3), and the corollary therefore

follows from Proposition 1 (Chapter IX, § 2, no. 4, Theorem 1, Corollary 2).

Note that this corollary applies in particular if X is locally compact and σ-compact (Chapter I, § 9, no. 9, Proposition 15, Corollary 1).

Now let Y be a metric space and let d be its metric. If X is any set and S any set of subsets of X, we shall denote by \(\mathcal{B}_{\mathfrak{S}}(X;Y)\) the set of all mappings \(u:X\to Y\) such that \(u(A)\) is bounded for each \(A\in S\) . Unless the contrary is expressly stated we shall regard \(\mathcal{B}_{\mathfrak{S}}(X;Y)\) as endowed with the uniformity of S-convergence, which is defined by the following family of pseudometrics on \(\mathcal{B}_{\mathfrak{S}}(X;Y)\) :

\[
d _ {\mathbf {A}} (u, v) = \sup _ {x \in \mathbf {A}} d (u (x), v (x))\tag{A ∈ S}
\]

which are finite by hypothesis. When \(\mathfrak{S} = \{\mathbf{X}\}\) , we write \(\mathcal{B}(\mathbf{X};\mathbf{Y})\) in place of \(\mathcal{B}_{\mathfrak{S}}(\mathbf{X};\mathbf{Y})\) . A mapping \(u:\mathbf{X}\to \mathbf{Y}\) is said to be bounded if it belongs to \(\mathcal{B}(\mathbf{X};\mathbf{Y})\) , i.e.~if \(u(\mathbf{X})\) is a bounded subset of \(\mathbf{Y}\) .

PROPOSITION 2. Let X be a set and Y a metric space. The set \(\mathcal{B}(X;Y)\) of bounded mappings is both open and closed in the space \(\mathcal{F}_{u}(X;Y)\) .

If \(u\) is bounded, then every mapping \(v: \mathbf{X} \to \mathbf{Y}\) such that for all \(x \in \mathbf{X}\) , we have \(d(u(x), v(x)) \leqslant 1\) is bounded, because

\[
d (v (x), v (x _ {0})) \leqslant d (u (x), u (x _ {0})) + 2;
\]

hence \(\mathcal{B}(\mathbf{X};\mathbf{Y})\) is open. On the other hand, if \(u\) lies in the closure of \(\mathcal{B}(\mathbf{X};\mathbf{Y})\) in \(\mathcal{F}_u(\mathbf{X};\mathbf{Y})\) , there is a mapping \(u_0\in \mathcal{B}(\mathbf{X};\mathbf{Y})\) such that \(d(u(x),u_0(x))\leqslant 1\) for all \(x\in \mathbf{X}\) ; hence \(u\) is bounded.

COROLLARY I. Let X be a set and Y a metric space. Then \(\mathcal{B}_{\mathfrak{S}}(X;Y)\) is closed in \(\mathcal{F}_{\mathfrak{S}}(X;Y)\) . In particular, if Y is complete then \(\mathcal{B}_{\mathfrak{S}}(X;Y)\) is complete with respect to the uniformity of \(\mathfrak{S}\) -convergence.

For \(\mathcal{B}_{\mathfrak{S}}(\mathbf{X};\mathbf{Y})\) is the inverse image of the subset \(\prod_{\mathbf{A}\in \mathfrak{S}}\mathcal{B}(\mathbf{A};\mathbf{Y})\) of the product \(\prod_{\mathbf{A}\in \mathfrak{S}}\mathcal{F}_u(\mathbf{X};\mathbf{Y})\) under the canonical mapping of \(\mathcal{F}_{\mathfrak{S}}(\mathbf{X};\mathbf{Y})\) into \(\prod_{\mathbf{A}\in \mathfrak{S}}\mathcal{F}_u(\mathbf{A};\mathbf{Y})\) ; the first assertion therefore follows from § 1, no. 2, Remark 3, and the second follows from the first, if we take account of Theorem 1 of § 1, no. 5.

COROLLARY 2. Let X be a topological space and Y a metric space. Then the space of all bounded continuous mappings of X into Y is both open and closed in \(\mathcal{C}_u(X;Y)\) ; it is complete if Y is complete.

The space in question is \(\mathcal{B}(\mathbf{X};\mathbf{Y})\cap \mathcal{C}_u(\mathbf{X};\mathbf{Y})\) ; the first assertion follows from Proposition 2; the second follows from the first (§ 1, no. 6, Theorem 2, Corollary 1).

